\section{Жадный алгоритм (GreedySolver)}

GreedySolver реализует простой, но эффективный жадный подход.

\subsection{Алгоритм}

\begin{enumerate}
    \item Коробки предварительно сортируются по убыванию площади основания
    \item Для каждой коробки в отсортированном порядке:
    \begin{enumerate}
        \item Генерируются все возможные ориентации коробки (до 6 вариантов)
        \item Для каждой ориентации получаем кандидатные позиции из HeightHandler
        \item Выбирается позиция и ориентация с минимальной итоговой высотой $h + h_{box}$
        \item Коробка размещается, HeightHandler обновляется
    \end{enumerate}
\end{enumerate}

\subsection{Сложность}

Временная сложность: $O(n \cdot k \cdot m)$, где $n$ --- число коробок, $k$ --- число ориентаций (до 6), $m$ --- число кандидатных позиций (растёт с числом уложенных коробок).

\subsection{Результаты}

На тестовом наборе из 436 задач GreedySolver показывает:
\begin{itemize}
    \item Среднее заполнение: 73.56\%
    \item Минимум: 59.01\%
    \item Максимум: 85.21\%
    \item Время: 308 мс (на все 436 тестов)
\end{itemize}

\section{Жадный поиск в большой окрестности (LNSSolver)}

LNSSolver использует метаэвристический подход Large Neighborhood Search (LNS) --- поиск в большой окрестности с жадным критерием принятия решений.

\subsection{Идея метода}

В отличие от классического локального поиска, который рассматривает небольшие изменения решения, LNS работает с ``большими окрестностями'' --- существенными модификациями текущего решения. Это позволяет эффективнее исследовать пространство решений и избегать локальных минимумов.

\subsection{Представление решения}

Решение кодируется как массив структур \texttt{BoxMeta}:
\begin{itemize}
    \item \texttt{box\_id} --- индекс коробки в исходном наборе
    \item \texttt{rotation} --- выбранная ориентация ($-1$ = рассмотреть все, $0$--$5$ = конкретная ориентация)
    \item \texttt{position} --- предпочтительный угол паллеты ($0$--$3$)
\end{itemize}

\textbf{Порядок элементов в массиве определяет порядок укладки коробок} --- это ключевая особенность кодирования. Перестановка элементов массива приводит к совершенно другой укладке.

\subsection{Функция симуляции}

Функция \texttt{simulate} принимает массив BoxMeta и симулирует укладку, используя жадный алгоритм:
\begin{enumerate}
    \item Для каждой коробки в заданном порядке:
    \begin{enumerate}
        \item Если \texttt{rotation} $\neq -1$ --- используем указанную ориентацию, иначе перебираем все допустимые ориентации
        \item Для каждой ориентации получаем кандидатные позиции из HeightHandler
        \item Выбираем позицию с минимальной итоговой высотой $h + h_{box}$
        \item При равной высоте предпочитаем позицию ближе к указанному углу (\texttt{position})
    \end{enumerate}
    \item Возвращается итоговая высота паллеты и метрики укладки
\end{enumerate}

\subsection{Операторы мутации}

Реализовано 13 операторов мутации, разделённых на категории:

\textbf{Изменение параметров (position/rotation):}
\begin{itemize}
    \item \texttt{position\_rotation\_fixed} --- установка одинакового значения для сегмента
    \item \texttt{position\_rotation\_random} --- случайные значения для каждой коробки в сегменте
\end{itemize}

\textbf{Изменение порядка:}
\begin{itemize}
    \item \texttt{reverse} --- реверс подмассива
    \item \texttt{shuffle} --- случайная перестановка подмассива
    \item \texttt{swap} --- обмен двух случайных коробок
    \item \texttt{swap\_adjacent} --- обмен соседних коробок
\end{itemize}

\textbf{Сортировки:}
\begin{itemize}
    \item \texttt{sort\_by\_area} --- сортировка по площади основания
    \item \texttt{sort\_by\_volume} --- сортировка по объёму
    \item \texttt{sort\_by\_height} --- сортировка по высоте
    \item \texttt{sort\_by\_weight} --- сортировка по весу
\end{itemize}

\textbf{Специальные:}
\begin{itemize}
    \item \texttt{threshold} --- изменение порога минимальной опоры
    \item \texttt{group\_by\_sku} --- группировка коробок одного типа вместе
    \item \texttt{move\_bad\_box} --- перемещение коробки с худшей опорой в начало
\end{itemize}

Каждому оператору назначается вес, определяющий вероятность его выбора.

\subsection{Критерий принятия}

Алгоритм использует \textbf{жадный критерий принятия}: новое решение принимается только если оно строго лучше текущего лучшего:

$$score_{new} > score_{best}$$

где $score = percolation + 0.3 \cdot min\_support\_ratio$

Это означает, что алгоритм никогда не ухудшает текущее лучшее решение, а только улучшает его. Такой подход проще имитации отжига и на практике показывает хорошие результаты за счёт большого разнообразия операторов мутации.

\subsection{Результаты}

\begin{table}[ht]
\centering
\caption{Результаты LNSSolver при разном времени работы}
\begin{tabular}{|l|c|c|c|c|}
\hline
\textbf{Время} & \textbf{Среднее} & \textbf{Мин} & \textbf{Макс} & \textbf{Общее время} \\
\hline
1 сек & 77.18\% & 72.86\% & 86.39\% & 14 сек \\
5 сек & 78.23\% & 73.11\% & 88.15\% & 70 сек \\
30 сек & 78.93\% & 75.36\% & 88.65\% & 7 мин \\
300 сек & 79.71\% & 76.03\% & 90.39\% & 70 мин \\
\hline
\end{tabular}
\end{table}

% TODO: Добавить график зависимости percolation от времени (figures/lns_convergence.png)
% Показать: ось X - время (логарифмическая шкала), ось Y - среднее percolation (%)
% Видна логарифмическая зависимость: основной прирост в первые секунды
